% ---- pg-270 （书页 258） ----

根时, 它的功效才最佳, 同时也是最大. 这接近于下面这种情形: 较大的差别恰好存在于某一特征指标中, 即存在于一个组与其余各组之间. 这里还有某些迹象表明, 皮莱迹检验对非正态性的适应能力要比其余方法稍强些. 但是, 当残差具有非正态性时, 我们还是建议先对原始数据作变换.

所有这四个检验统计量均可应用于双因子情形, 甚至应用于更为复杂的 MANOVA 中, 详细的讨论将在第 7 章的多元回归模型中给出.

当且仅当多元检验表明有差异或者同零假设有偏离时, 我们才较深入地探查其原因. 我们建议对所有组对及所有特征计算庞弗罗尼区间. 典型的情况是: 由置信椭球的投影确定的联合置信区间太大. 单一区间也许对值得进一步研究的那些差异有启发性, 但就现有数据而论, 尚不能将这些区间当成作出“存在差异”这种结论的证据. 我们就本章提出的比较各种处理方法作一小结. 第一步是利用目测法和计算方法检查一下数据, 看看是否存在离群值.

\begin{mvframe}[对各种处理进行多元比较的策略]
\begin{enumerate}
\item 识别离群值. 对数据进行分批检查, 看有无离群值, 同时对来自任何经过拟合的模型的所有残差向量进行检查, 看有无离群值.
\item 对假设完成多元检验. 我们选择似然比检验, 它等价于威尔克斯 $\Lambda$ 检验.
\item 计算庞弗罗尼联合置信区间. 如果多元检验揭示存在某种差异, 则对所有的组对或处理以及所有的特征指标, 计算它们的庞弗罗尼置信区间. 如果所有差异均不显著, 那就对包含差异的较大的响应集与响应对之和求庞弗罗尼区间.
\end{enumerate}
\end{mvframe}

对上面提出的策略必须注意一点, 即差异可能仅出现在许多特征指标中的某一个之中. 此外, 差异仅出现在很少的一些处理组合之中. 于是, 这少数起作用的差异有可能消失在所有不起作用的差异之中. 也就是说, 整体检验可能显示不显著, 而在对某一特定的起作用变量作一元检验时, 就会发现存在差异. 避免出现这种情况的方法是要构造一个好的实验设计. 当一个特定变量预期会产生差异时, 就设计一个有效的实验, 使它不把太多的预期不会显示诸处理之间差异的其他变量包括进去.

\vspace{6pt}
{\large\bfseries 练习}

\medskip
\noindent 6.1\quad 利用例 6.1 中的污水数据及结论, 构造并画出均值向量差 $\boldsymbol{\delta}$ 的 95\% 联合置信区域. 注意到 $\boldsymbol{\delta}=\mathbf{0}$ 这个点落在 95\% 置信区域之外, 试问这一结果是否与例 6.1 中考虑的假设 $H_0: \boldsymbol{\delta}=\mathbf{0}$ 的检验结果一致? 请作解释.

\noindent 6.2\quad 利用例 6.1 提供的信息, 构造均值差向量 $\boldsymbol{\delta}$ 诸分量的 95\% 庞弗罗尼联合置信区间, 并将这些区间的宽度与例 6.1 中所给的联合置信区间的宽度作比较.

\noindent 6.3\quad 与\hyperref[tab:6.1]{表 6.1} 中的样本 8 相对应的数据似乎大得异常. 摒弃样本 8, 构造均值差向量 $\boldsymbol{\delta}$ 的 95\% 联合置信域以及该向量诸分量的 95\% 庞弗罗尼联合置信区间, 结论是否与 $H_0: \boldsymbol{\delta}=\mathbf{0}$ 的检验相一致? 试作讨论. 这一离群数据是否会使这些数据的分析产生差异?

\noindent 6.4\quad 参照例 6.1:

\noindent (a) 对例 6.1 中各对观测值作对数变换后, 重新进行分析.

\noindent (b) 对变换后的变量, 构造其均值差向量 $\boldsymbol{\delta}$ 诸分量的 95\% 庞弗罗尼联合置信区间.

\noindent (c) 对变换后数据的差值向量而言, 关于二元正态分布这一假定是否有可能遭到破坏?

\noindent 6.5\quad 一名研究人员用三个指标来衡量心脏病发作的严重程度, 从某医院急诊室的 $n=40$ 名心脏病患者测得这些指标, 其样本均值向量与协方差阵分别为
